\begin{exo}{}
	Dans le repère orthonormé $\vOIJ$ ci-dessous, on a placé trois points $R$, $S$ et $T$.

	\ms

	\begin{center}
		\begin{tikzpicture}[scale=0.5, >=Stealth]
			\draw[gray!30, thick] (-0.5,-0.5) grid (14.5,5.5);
			\draw[thick, ->] (-0.5,0) -- (14.5,0);
			\draw[thick, ->] (0,-0.5) -- (0,5.5);
			\foreach \x in {1,2,3,4,5,6,7,8,9,10,11,12,13,14} { \draw (\x,0.1) -- (\x,-0.1); }
			\foreach \y in {1,2,3,4,5} { \draw (0.1,\y) -- (-0.1,\y); }
			\draw[thick, ->] (0,0) -- (1,0) node[midway, below] {$\vec i$};
			\draw[thick, ->] (0,0) -- (0,1) node[midway, left] {$\vec j$};
			\node[below left] at (0,0) {$O$};
			\coordinate (R) at (2,1);
			\coordinate (S) at (9,3);
			\coordinate (T) at (13,4);
			\node[above right, red!80!black] at (R) {$R$};
			\node[above right, red!80!black] at (S) {$S$};
			\node[above right, red!80!black] at (T) {$T$};
			\foreach \p in {R,S,T} {
					\draw[thick, red!80!black] ($(\p)+(-0.15,-0.15)$) -- ($(\p)+(0.15,0.15)$);
					\draw[thick, red!80!black] ($(\p)+(-0.15,0.15)$) -- ($(\p)+(0.15,-0.15)$);
				}
		\end{tikzpicture}
	\end{center}

	\begin{enumerate}
		\item Lire les coordonnées des points $R$, $S$ et $T$.
		\item Calculer la valeur exacte de $RS$, $ST$ et $RT$.
		\item À l'aide de la calculatrice, donner une valeur approchée au millième de chacune de ces distances.
		\item L'égalité $RS + ST = RT$ est-elle vérifiée ?
		\item Que peut-on en conclure concernant l'alignement des points $R$, $S$ et $T$ ?
	\end{enumerate}
\end{exo}
